\documentclass[10pt,a4paper]{amsart}
\usepackage[margin=19mm]{geometry}
\usepackage{amsmath,amssymb,mathtools}
\usepackage[scr=rsfs,cal=euler]{mathalfa}
\usepackage{xfrac}
\usepackage[unicode,hidelinks,bookmarks=false]{hyperref}
\newcommand{\ie}{i.e.\ }
\newcommand{\cf}{cf.\ }
\newcommand{\resp}{respectively\ }
\newcommand{\Subsection}{\subsection}
\providecommand{\nobreakdash}{\nobreak\hbox{-}}
\setlength{\emergencystretch}{2em}
\multlinegap0pt
\usepackage{amssymb}
\usepackage{tikz}
\usepackage{float}
\newtheorem{definition}{Definition}
\newtheorem{proposition}{Proposition}
\newcommand{\B}{\mathcal B}
\newcommand{\F}{\mathcal F}
\title{B\'ezout coefficients of coprime numbers approximate quadratic B\'ezier curves}
\author{\sc{Benjam\'in A.~Itz\'a-Ortiz}, \sc{Roberto L\'opez-Hern\'andez}, \sc{Pedro Miramontes}}
\date{Source: arXiv:2012.11770v2 (CC BY 4.0)}
\begin{document}
\maketitle

\noindent\emph{Source and licence.} B\'ezout coefficients of coprime numbers approximate quadratic B\'ezier curves. Version arXiv:2012.11770v2 (CC BY 4.0), distributed by arXiv under the Creative Commons Attribution 4.0 International licence, http://creativecommons.org/licenses/by/4.0/, as recorded in the arXiv licence metadata for this exact version and shown on the abstract page https://arxiv.org/abs/2012.11770v2. Full licence notice: \texttt{licenses/arxiv-2012.11770-CC-BY-4.0.txt}.\par\smallskip

\noindent\textit{Editorial scope.} Counted unit: the article's proposition mainprop (source0.tex lines 169--180). The article's complete original proof of it is delivered with it (source0.tex lines 181--237, one \texttt{proof} environment of 57 lines). No mathematics is written, repaired or re-derived: the statement and the proof are the article's own lines, in the article's own notation, and no retained line is substituted or re-flowed.  Retained uncounted as the source-local context the proof needs: lines 116--123; lines 129--135.  Nothing was omitted from any retained span, so the package carries no omission record.  No retained line contains a citation, so the wrapper carries no bibliography and no bibliography entry.  No retained line contains a figure environment, an image-inclusion command or drawing code, so no image is delivered and the package declares no asset dependency.  Editorial formatting only, in addition: an AMS \texttt{amsart} preamble that restates the article's own declarations and macros used by the retained lines, namely the environments \texttt{definition}, \texttt{proposition} on separate counters, together with the standard packages those lines need.  The retained environments print in this document as Definition 1, Proposition 1 in order of appearance, whereas the article prints the same items with the numbering of its own full document.  The complete native source of the article as distributed in the arXiv e-print for this exact version is bundled unchanged as \texttt{sources/106136/source0.tex}.
\begin{definition}\label{defbezout}
Given a pair of positive coprime integers $(p,q)$, we define the {\em B\'ezout coefficients} of $(p,q)$ 
as the unique pair of coprime numbers, denoted $\mathcal B(p,q)=(a,b)$, such that $0<a \leq p$, $0\leq b < q$  and  satisfy  {B\'ezout's identity}
\begin{equation*}
    aq-bp=1.
\end{equation*}

 \end{definition}

\begin{proposition}\label{prop}
Let $(p,q)$ be a coprime pair of nonnegative integers. If $\B(p,q)=(a,b)$ then following identities hold.
    \begin{enumerate} 
    \item  $\B(q,p)=(q-b,p-a)$.
    \item  $\B(b+q,a+p)=(q,p)$.
    \end{enumerate}
\end{proposition}

\begin{proposition}\label{mainprop}
Let $p$ and $q$ be relatively prime positive numbers. The following conditions hold.
   \begin{enumerate}
    \item  $\|\B(p,q)-(p,q)\|=\|\B(q,p)-(0,0)\|$.
    %   \item $\B\F\B(q,p)=\B(p,q)
    %  \Rightarrow \B\F\B(p,q)\not=\B(q,p)$.
    \item The distance from the point $\B(p,q)$ to the line $y=\dfrac{q}{p}x$  and the distance from the point $\B(q,p)$ to the line $y=\dfrac{p}{q}x$ are both  equal to $\dfrac{1}{\sqrt{p^2+q^2}}$. 
    \item Let $t_0=1-\dfrac{\B(p,q)\cdot (p,q)}{\|(p,q)\|^2}$. Then $0<t_0<1$ and the projection of the point $\B(p,q)$ on the line $y=\dfrac{q}{p}x$ is  $(1-t_0)(p,q)$, while the projection of the point $\B(q,p)$ on the line $y=\dfrac{p}{q}x$ is $t_0(q,p)$. 
    \item The distance from the point $(p,q)$ to the projection of $\B(p,q)$ on the line $y=\dfrac{q}{p}x$ is
    equal to the distance from the origin $(0,0)$ to the projection of $\B(q,p)$  on $y=\dfrac{p}{q}x$.
    \end{enumerate}
\end{proposition}

\begin{proof}
   
Let $p$ and $q$  be relatively prime positive integers. Assume $\B(p,q)=(a,b)$. From part (1) in Proposition~\ref{prop} we see that $\B(q,p)=(q-b,p-a)$. Hence
\begin{align*}
    \|\B(q,p) -(0,0)\| &= \sqrt{(q-b)^2+(p-a)^2}\\
                       &=\| (q,p)-(b,a)\|\\
                       &= \| (a,b)-(p,q)\|\\
                       &=\| \B(p,q)-(p,q)\|,
\end{align*}
thus proving (1).


For $(2)$, we compute the distance from $\B(p,q)=(a,b)$ to $y=\dfrac{q}{p}x$ to be 
\begin{align*}
    \dfrac{\left| aq-bp\right|}{\sqrt{p^2+q^2}} &= \dfrac{1}{\sqrt{p^2+q^2}}.
\end{align*}
 Similarly, if $\B(p,q)=(a,b)$ then the distance from $\B(q,p)=(q-b,p-a)$ to $y=\dfrac{p}{q}x$ is
\begin{align*}
    \dfrac{\left| p(q-b)-q(p-a)\right|}{\sqrt{p^2+q^2}} &= \dfrac{1}{\sqrt{p^2+q^2}}.
\end{align*}

Now, to prove (3), let us assume again that $\B(p,q)=(a,b)$.  Then, by Definition~\ref{defbezout}, we have the inequalities $0\leq a <p$ and $0<b\leq q$. %Then $cr-dp=1$ and so $cq-dp=1-hc$.
We obtain that the vector resolution of $\B(p,q)$ in the direction of $(p,q)$ has length less than $\|(p,q)\|$, that is, $0< \dfrac{\B(p,q)\cdot (p,q)}{\|(p,q)\|}< \|(p,q)\|$. Thus $0< t_0 =1-\dfrac{\B(r,s)\cdot (r,s)}{\|(r,s)\|^2}< 1$. 

By Proposition~\ref{prop} we then have $\B(q,p)=(q-b,p-a)$. The coordinates of the projection of   $\B(p,q)$ on $y=\dfrac{q}{p}x$ are nothing but
\begin{equation*}
    \dfrac{\B(p,q)\cdot (p,q)}{\|(p,q)\|^2} (p,q)=(1-t_0)(p,q),
\end{equation*}
while the coordinates of the projection of $B(q,p)$ on $y=\dfrac{p}{q}x$ are
\begin{equation*}
    \dfrac{\B(q,p)\cdot (q,p)}{\|(q,p)\|^2} (q,p).
\end{equation*}
 By computing
 \begin{align*}
     t_0&=1-\dfrac{\B(p,q)\cdot (p,q)}{\|(p,q)\|^2} \\
     &=\dfrac{(p,q)\cdot (p,q)}{\|(p,q)\|^2}-\dfrac{(a,b)\cdot (p,q)}{\|(p,q)\|^2}\\
     &=\dfrac{(p-a,q-b)\cdot (p,q)}{\|(p,q)\|^2}\\
     &=\dfrac{\B(q,p)\cdot (q,p)}{\|(p,q)\|^2}\\
\end{align*}
 we complete the proof of $(3)$.
 
 
 
  Finally, for the proof of $(4)$, notice that the triangle with vertices $\B(p,q)$, $(p,q)$ and the point on the line $qx-py=0$ closest to $\B(p,q)$ is congruent to the  triangle with vertices $\B(q,p)$, $(0,0)$ and the point on the line $px-qy=0$ closest to $\B(q,p)$, by parts (1) and (2). Thus, (4) follows.  We may also obtain the result by direct computation using part $(3)$  as follows.
\begin{align*}
    \left\|\dfrac{\B(q,p)\cdot (q,p)}{\|(q,p)\|^2} (q,p) -(0,0)  \right\|
    &= \left\| t_0 (q,p)   \right\|  \\
    &=\left\| t_0 (p,q)   \right\|  \\
    &=\left\|
    \left(1-\dfrac{\B(p,q)\cdot (p,q)}{\|(p,q)\|^2}\right) (p,q)
    \right\|\\
    &=\left\|(p,q)-
    \dfrac{\B(p,q)\cdot (p,q)}{\|(p,q)\|^2} (p,q)
    \right\|,
\end{align*}
as was to be proved.
\end{proof}

\end{document}
